% Generated by build_manuscript.py from 02_related_work.md
\section{2 Related work}\label{related-work}

\subsection{2.1 Image-based residual-strength assessment}\label{image-based-residual-strength-assessment}

Image-based assessment connects observable damage with a mechanical quantity that matters after impact. The public datasets of Hasebe and colleagues provide complementary resources: post-impact surface and internal images, and separately recorded compression-after-impact measurements \citep{hasebe2022data, hasebe2025data}. These resources support specimen-level correspondence between imaging and strength while preserving the difference between a damage observation and a load-bearing property. That distinction motivates the use of CAI error as the acquisition objective here. A region that is visually conspicuous is a candidate for inspection, but conspicuity alone does not establish its value for the strength estimate.

A direct journal neighbour is Mack et al.'s ResNet18-based framework for predicting impact energy and CAI strength from C-scan damage images \citep{mack2026}. The authors' institutional abstract describes direct image interpretation and an analysis of feature-importance maps. This establishes image-to-strength regression as an existing research direction. The distinction in the present study is the sequential allocation of partial observations and the associated quality--acquisition trajectory. We do not compare their reported R² numerically with ours because the data preparation and evaluation protocols differ.

Image regression and selective observation address related but separable questions. A predictor must interpret whatever information is supplied, while an acquisition policy must decide which missing information to request. Improving either component can improve the combined system, but changing both together complicates interpretation of a comparison. The common frozen predictor used here deliberately fixes the assessment mapping across acquisition strategies. This design makes the supplied subsets and their ordering the experimental focus, rather than claiming that the selected regressor is the strongest possible model for complete C-scan images.

\subsection{2.2 Adaptive acquisition for nondestructive inspection}\label{adaptive-acquisition-for-nondestructive-inspection}

Autonomous ultrasonic acquisition already has a substantial conceptual precedent. Fuentes et al.~formulate the selection of inspection locations through Bayesian optimization and robust outlier analysis \citep{fuentes2020}. Their institutional abstract describes sequentially updating a two-dimensional field of novelty indices and selecting locations with evidence of damage, with component damage probability as an output. Thus, choosing the next ultrasonic observation from information acquired so far is not introduced by the present work. The pertinent difference is the task endpoint: we optimize a CAI regression trajectory under an image-pixel budget rather than the reported damage-indication objective.

This distinction also clarifies the status of the fixed controls. A geometry-spread order measures broad spatial coverage, while center-first concentrates observations near the image centre. Neither reproduces an entire autonomous inspection system, and serpentine image order alone does not model robotic travel efficiency. The controls instead provide transparent alternatives under identical replay rules. The literature motivates adaptive inspection; the empirical comparison in this paper concerns the specified nine acquisition strategies and their shared strength predictor.

\subsection{2.3 Task-driven sequential information acquisition}\label{task-driven-sequential-information-acquisition}

Dynamic feature acquisition provides the broader learning formulation. Shim et al.~jointly learn a classifier and an acquisition agent, using a set representation of observed features and actions that either acquire information or stop and predict \citep{shim2018}. Janisch et al.~formulate costly-feature classification as sequential decision making, with feature requests and classification actions \citep{janisch2019}. Their formulations show that prediction quality and information cost can be optimized together. The present policy instead uses a fixed acquisition cap and a frozen regression model, without a learned classification or stopping action. This is an application-specific choice that preserves a common assessment model for the comparison, rather than a general argument against joint learning.

Covert et al.~develop a greedy dynamic-selection method based on conditional mutual information and amortized optimization \citep{covert2023}. Their formulation uses a fixed feature-count budget with uniform costs; its regression result relates squared-error optimality to conditional variance. Our whole-cell native-pixel costs and absolute-error trajectory objective differ from those assumptions. We therefore use this work to position task-relevant acquisition, not as an optimality guarantee or a reproduced baseline. Policy-gradient learning supplies the optimization mechanism in our implementation \citep{williams1992}, while the manuscript contribution lies in the inspection-state contract, common-predictor comparison and process-level evidence.

Together, these lines locate the work between image-based mechanical assessment and sequential information acquisition. The method combines a surface-derived starting prior with evidence-dependent internal acquisition, then evaluates both intermediate quality and the final strength estimate. This connection is useful precisely because damage indication, classification utility and CAI regression error need not assign the same value to an observation. The closest-work matrix in the supplementary material records the verified comparison dimensions and leaves inaccessible details explicitly unverified.
